\documentclass[11pt]{article}
\usepackage[a4paper,margin=0.65in]{geometry}
\usepackage[T1]{fontenc}
\usepackage{amsmath,amssymb,amsfonts}
\usepackage{graphicx}
\usepackage{pdfpages}
\usepackage[english,shorthands=off]{babel}
\usepackage{fancyhdr}
\usepackage[bookmarks=false,hidelinks]{hyperref}
\usepackage[protrusion=true,expansion=false]{microtype}
\emergencystretch=3em
\tolerance=600
\providecommand{\modd}{\mathrm{mod}}
\providecommand{\Disc}{\mathrm{Disc}}
\providecommand{\canont}{}
\providecommand{\asterism}{*}
\providecommand{\Hessian}{H}
\providecommand{\tome}{}
\providecommand{\page}{p.}
\pagestyle{fancy}\fancyhf{}\fancyhead[L]{Pages 251--262}\fancyhead[R]{Cayley --- Collected Papers, Vol.~I}\renewcommand{\headrulewidth}{0.4pt}
\setlength{\headheight}{15pt}

\begin{document}

\begin{center}
{\LARGE\bfseries The Collected Mathematical Papers of}\\[0.5em]
{\LARGE\bfseries Arthur Cayley}\\[1em]
{\Large Volume I}\\[1em]
{\large PDF Pages 251--262 \quad (Book pages 229--240)}\\[2em]
\end{center}

\tableofcontents
\newpage

% ===========================================================================
% PDF page 251 -- Book page 229 -- Paper [34] continued
% ===========================================================================
\section*{34.}
\section*{NOTE ON THE MAXIMA AND MINIMA OF FUNCTIONS OF THREE VARIABLES.}
\addcontentsline{toc}{section}{34. Note on the Maxima and Minima of Functions of Three Variables}
\setcounter{page}{229}

\noindent\textit{[Continuation of paper 34.]}

\bigskip

It is required to show that if $A' + B' + C'$ and $A'' + B'' + C''$ are positive,
$A'$, $B'$, $C'$, $A''$, $B''$, $C''$ are so likewise.

Consider the cubic equation
\[
(A'-k)(B'-k)(C'-k)-(A'-k)F^2-(B'-k)G^2-(C'-k)H^2+2FGH=0,
\]
the roots of which are all real. By the formulae just given this may be written
\[
k^3-k^2(A'+B'+C')+k(A''+B''+C'')-K^2=0;
\]
and the terms of this equation are alternately positive and negative; i.e. the roots are
all positive. Hence the roots of the limiting equation
\[
(B'-k)(C'-k)-F^2=0
\]
are positive, i.e. $B'+C'$ and $B'C'$ are positive: but from the second condition
$B'$, $C'$ are of the same sign: consequently they are of the same sign with
$B'+C'$, or positive. Also $A''=B'C'-F^2$ is positive. Similarly, considering the
other limiting equations, $A'$, $B'$, $C'$, $A''$, $B''$, $C''$ are all of them positive.

In connection with the above I may notice the following theorem. The roots of the
equation
\[
\begin{aligned}
&(A-ka)(B-kb)(C-ck)-(A-ka)(F-kf)^2-(B-kb)(G-kg)^2\\
&\qquad{}-(C-kc)(H-kh)^2+2(F-kf)(G-kg)(H-kh)=0,
\end{aligned}
\]
are all of them real, if either of the functions
\[
Ax^2+By^2+Cz^2+2Fyz+2Gzx+2Hxy,
\]
\[
ax^2+by^2+cz^2+2fyz+2gxz+2hxy,
\]
preserve constantly the same sign. The above form parts of a general system of
properties of functions of the second order.

\newpage
% ===========================================================================
% PDF page 252 -- Book page 230 -- Paper [35]
% ===========================================================================
\section*{35.}
\section*{ON HOMOGENEOUS FUNCTIONS OF THE THIRD ORDER WITH THREE VARIABLES.}
\addcontentsline{toc}{section}{35. On Homogeneous Functions of the Third Order with Three Variables}
\setcounter{page}{230}

\textit{[From the Cambridge and Dublin Mathematical Journal, vol.~I. (1846), pp.~97--104.]}

\bigskip

The following problem corresponds to the geometrical question of determining the
polar reciprocal of a plane curve of the third order: the solution of it is also important,
with reference to the linear transformations of homogeneous functions of three variables
of the third order; reasons for which it has appeared to me worth while to obtain the
completely developed result.

Let
\[
\begin{aligned}
3U={}&ax^3+by^3+cz^3+3iy^2z+3jz^2x+3kx^2y+3i_1yz^2\\
&\quad{}+3j_1zx^2+3k_1xy^2+6lxyz .
\end{aligned}\tag{1}
\]
It is required to eliminate $x$, $y$, $z$, $\lambda$ from the equations
\[
U=0,\tag{2}
\]
\[
\left.
\begin{aligned}
\frac{dU}{dx}+\lambda\xi&=0,\\
\frac{dU}{dy}+\lambda\eta&=0,\\
\frac{dU}{dz}+\lambda\zeta&=0,
\end{aligned}
\right\}\tag{3}
\]
From the equations (2), (3), we obtain immediately
\[
\Theta=\xi x+\eta y+\zeta z=0,\tag{4}
\]
and thence
\[
\Theta x=0,\qquad \Theta y=0,\qquad \Theta z=0.\tag{5}
\]

\newpage
% ===========================================================================
% PDF page 253 -- Book page 231
% ===========================================================================
\setcounter{page}{231}

so that a single equation more, such as
\[
\Phi=0,\tag{6}
\]
where $\Phi$ is homogeneous and of the second order in $x$, $y$, $z$, would, in
conjunction with the equations (3) and (5), enable us to eliminate linearly the seven
quantities $x^2$, $y^2$, $z^2$, $yz$, $zx$, $xy$, $\lambda$. Such an equation may be
thus obtained.

Let $L$, $M$, $N$, $R$, $S$, $T$, be the second differential coefficients of $U$, each of
them divided by two. The equations (3) may be written
\[
\left.
\begin{aligned}
Lx+Ty+Sz+\lambda\xi&=0,\\
Tx+My+Rz+\lambda\eta&=0,\\
Sx+Ry+Nz+\lambda\zeta&=0.
\end{aligned}
\right\}\tag{7}
\]
And joining to these the equation (4),
\[
\xi x+\eta y+\zeta z=0,
\]
we have, by the elimination of $x$, $y$, $z$, in so far as they explicitly appear, and
$\lambda$, an equation $\Phi=0$ of the required form. Hence we may write
\[
\Phi=-\left|
\begin{array}{cccc}
L&T&S&\xi\\
T&M&R&\eta\\
S&R&N&\zeta\\
\xi&\eta&\zeta&0
\end{array}
\right|,\tag{8}
\]
or substituting for $L$, $M$, $N$, $R$, $S$, $T$, and expanding,
\[
\Phi=Ax^2+By^2+Cz^2+2Fyz+2Gzx+2Hxy,\tag{9}
\]
where the values of $A$, $B$, $C$, $F$, $G$, $H$ are \dotfill (10).

[I omit these values (10), and the values in the subsequent equations (13), (14),
(20): the values (10) and (13) serve for the calculation of (14), $FU$, the expression
for which with the letters $(a,b,c,f,g,h,i,j,k,l)$ in place of
$(a,b,c,i,j,k,i_1,j_1,k_1,l)$ is reproduced in my ``Third Memoir on Quantics,''
\textit{Phil. Trans.} vol.~CXLVI. (1856) pp.~627--647: the values (20) serve for finding
that of (21), $K(U)$, but the developed expression of this has not been calculated.]

Performing the elimination indicated, the result may be represented by
\[
FU=\left|
\begin{array}{ccccccc}
a&k_1&j&l&j_1&k&\xi\\
k&b&i_1&i&l&k_1&\eta\\
j_1&i&c&i_1&j&l&\zeta\\
2\xi&0&0&0&\zeta&\eta&0\\
0&2\eta&0&\zeta&0&\xi&0\\
0&0&2\zeta&\eta&\xi&0&0\\
A&B&C&F&G&H&0
\end{array}
\right|=0.\tag{11}
\]
Partially expanding,
\[
FU=Aa+Bb+Cc+2Ff+2Gg+2Hh.\tag{12}
\]

\newpage
% ===========================================================================
% PDF page 254 -- Book page 232
% ===========================================================================
\setcounter{page}{232}

The values of the coefficients $a$, $b$, $c$, $f$, $g$, $h$ may be useful on other
occasions: they are as follows \dotfill (13).

Substituting these values the result after all reductions becomes
\[
0=F(U)=\dotfill\tag{14}
\]
It would be desirable, in conjunction with the above, to obtain the equation
\[
K(U)=0,
\]
which results from the elimination of $x$, $y$, $z$ from the equations
\[
\frac{dU}{dx}=0,\qquad \frac{dU}{dy}=0,\qquad \frac{dU}{dz}=0,\tag{15}
\]
(i.e. the condition of a curve of the third order having a multiple point), but to effect
this would be exceedingly laborious. The following is the process of the elimination as
given by Dr Hesse, \textit{Crelle}, t.~xxviii. (and which applies also to the case of any
three equations of the second order). Forming the function $\nabla U$, of the third order
in $x$, $y$, $z$, by means of the equation
\[
\nabla U=\left|
\begin{array}{ccc}
L&T&S\\
T&M&R\\
S&R&N
\end{array}
\right|,\tag{16}
\]
($L$, $M$, $N$, $R$, $S$, $T$, the same as before).

Then, in consequence of the equations (15), we have not only
\[
\nabla U=0,\tag{17}
\]
which is very easily proved to be the case, but also
\[
\frac{d}{dx}\nabla U=0,\qquad
\frac{d}{dy}\nabla U=0,\qquad
\frac{d}{dz}\nabla U=0,\tag{18}
\]
as will be shown in a subsequent paper ``On Points of Inflection.'' [I think never
written.]

And from the six equations (15), (18), the six quantities $x^2$, $y^2$, $z^2$, $yz$, $zx$,
$xy$, may be linearly eliminated: we have
\[
\begin{aligned}
\nabla U={}&Ax^3+By^3+Cz^3+3Iy^2z+3Jz^2x+3Kx^2y+3I_1yz^2\\
&\quad{}+3J_1zx^2+3K_1xy^2+6\Lambda xyz,
\end{aligned}\tag{19}
\]
where the values of the coefficients $A$, $B$, \ldots $\Lambda$ are \dotfill (20),
and the result of the elimination is
\[
K(U)=\left|
\begin{array}{cccccc}
a&k_1&j&l&j_1&k\\
k&b&i_1&i&l&k_1\\
j_1&i&c&i_1&j&l\\
A&K_1&J&L&J_1&K\\
K&B&I_1&I&L&K_1\\
J_1&I&C&I_1&J&L
\end{array}
\right|=0.\tag{21}
\]
$[K(U)$ is consequently, as is well known, a function of the twelfth order in
$a$, $b$, $c$, $i$, $j$, $k$, $i_1$, $j_1$, $k_1$, $l]$.

\newpage
% ===========================================================================
% PDF page 255 -- Book page 233
% ===========================================================================
\setcounter{page}{233}

The equation
\[
\nabla U=0,
\]
combined with that of the curve, determine, as Dr Hesse has demonstrated in the paper
quoted, the points of inflection of the curve. It may be inferred from this, that if $U$
reduce itself to the form
\[
U=(\alpha x^2+\beta y^2+\gamma z^2+2\iota yz+2\kappa xz+2\lambda xy)P=VP,\tag{22}
\]
$P$ a linear function of $x$, $y$, $z$: then $\nabla U$ takes the form
\[
\nabla U=P(\rho V+\sigma P^2),\tag{23}
\]
where $\rho$ is of the second order in the coefficients of $P$, and also in the
coefficients $\alpha$, $\beta$, $\gamma$, $\iota$, $\kappa$, $\lambda$: and $\sigma$ is equal
to the determinant
\[
\left|
\begin{array}{ccc}
\alpha&\lambda&\kappa\\
\lambda&\beta&\iota\\
\kappa&\iota&\gamma
\end{array}
\right|,\tag{24}
\]
multiplied by a numerical factor. If $U$ is of the form
\[
U=PQR,\tag{25}
\]
then
\[
\nabla U=\rho PQR=\rho U,\tag{26}
\]
and this equation is consequently the condition of the function $U$ being resolvable
into linear factors. The equation in question resolves itself into
\[
\frac{A}{a}=\frac{B}{b}=\frac{C}{c}=\frac{I}{i}=\frac{J}{j}=\frac{K}{k}
=\frac{I_1}{i_1}=\frac{J_1}{j_1}=\frac{K_1}{k_1}=\frac{\Lambda}{l},\tag{27}
\]
a system which must contain three independent equations only. It would be interesting
to verify this \textit{a posteriori}.

\newpage
% ===========================================================================
% PDF page 256 -- Book page 234 -- Paper [36]
% ===========================================================================
\section*{36.}
\section*{ON THE GEOMETRICAL REPRESENTATION OF THE MOTION OF A SOLID BODY.}
\addcontentsline{toc}{section}{36. On the Geometrical Representation of the Motion of a Solid Body}
\setcounter{page}{234}

\textit{[From the Cambridge and Dublin Mathematical Journal, vol.~I. (1846), pp.~164--167.]}

\bigskip

Let $P$, $Q$, $R$, \ldots be consecutive generating lines of a skew surface, and on these
take points $p'$, $p$; $q'$, $q$; $r'$, $r$, \ldots such that $pq'$, $qr'$, \ldots are the
shortest distances between $P$ and $Q$, $Q$ and $R$, \&c. Then for the generating line
$P$, the ratio of the inclination of the lines $P$, $Q$ to the distance $pq'$ is said to be
``the torsion,'' the angle $q'pq$ is said to be the deviation, and the ratio of the
inclination of the planes $Qpq'$ and $Qqr'$ to the inclination of $P$ and $Q$ is said to
be the ``skew curvature.'' And similarly for any other generating line; so that the
torsion and deviation depend on the position of the consecutive line, and the skew
curvature on the position of the two consecutive lines. The curve $pqr\ldots$ is said to
be the minimum distance curve [or curve of striction]. \{When the skew surface
degenerates into a developable surface, the torsion is infinite, the deviation a right
angle, the skew curvature proportional to the curvature of the principal section, i.e. it
is the distance of a point from the edge of regression, multiplied into the reciprocal of
the radius of curvature, a product which is evidently constant along a generating line.
Also the curve of minimum distance becomes the edge of regression.\} A skew surface,
considered independently of its position in space, is determined when for each generating
line we know the torsion, deviation, and skew curvature. For, assuming arbitrarily the
line $P$ and the point $p$, also the plane in which $pq'$ lies, the position of $Q$ is
completely determined from the given torsion and deviation; and then $Q$ being known,
the position of $R$ is completely determined from the skew curvature for $P$, and the
torsion and deviation for $Q$; and similarly the consecutive generating lines are to be
determined.

Two skew surfaces are said to be ``deformations'' of each other, when for corresponding
generating lines the torsion is always the same. Thus a surface will be deformed if
considering the elements between the successive generating lines $P$, $Q$, \ldots as
rigid, these

\newpage
% ===========================================================================
% PDF page 257 -- Book page 235
% ===========================================================================
\setcounter{page}{235}

elements be made to revolve round the successive generating lines $P$, $Q$, \ldots and
to slide along them. \{They are ``transformations'', when not only the torsions but also
the deviations are equal at corresponding generating lines: thus, if the sliding of the
elements along $P$, $Q$, \ldots be omitted, the new surface will be, not a deformation,
but a transformation of the other.\} No two skew surfaces can be made to roll and slide
one upon the other, so that their successive generating lines coincide, unless one of
them is a deformation of the other: and when this is the case, the rolling and sliding
motions are \textit{completely determined}. In fact the angular velocity of the generating
line is the angular velocity round this line, into the difference of the skew curvatures
of the two surfaces; the velocity of translation of the generating line in its own
direction is to the angular velocity of the generating line, as the difference of the
deviations is to the torsion. \{This includes also the case in which one surface is a
transformation of the other, where the motion is evidently a rolling one.\} A skew
surface moving in this manner upon another of which it is the deformation, may be said
to ``glide'' upon it. We may now state the kinematical theorem:

\begin{quote}
``Any motion whatever of a solid body in space may be represented as the `gliding'
motion of one skew surface upon another fixed in space, and of which it is the
deformation.''
\end{quote}
a theorem which is to be considered as the generalization of the well-known one---
\begin{quote}
``Any motion of a solid body round a fixed point may be represented as the rolling
motion of a conical surface upon a second conical surface fixed in space.''
\end{quote}
and of the supplementary theorem---
\begin{quote}
``The angular velocity round the line of contact (the instantaneous axis) is to the
angular velocity of this line as the difference of curvatures of the two cones at any
point in the same line, to the reciprocal of the distance of the point from the vertex.''
\end{quote}

The analytical demonstration of this last theorem is rather interesting: it depends on
the following formulae. Forming two determinants, the first with the angular velocities
round three axes fixed in space, and the first and second derived coefficients of these
velocities with respect to the time; the other in the same way with the angular
velocities round axes fixed in the body; the difference of these determinants is equal
to the fourth power of the angular velocity into the square of the angular velocity of
the instantaneous axis.

To show this, let $p$, $q$, $r$ be the angular velocities round the axes fixed in the
body; $u$, $v$, $w$ those round axes fixed in space; $\omega$ the angular velocity round
the instantaneous axis; $\nabla$, $\Omega$ the two determinants: the theorem comes to
\[
\nabla-\Omega=M,
\]
where
\[
M=\omega^2(p'^2+q'^2+r'^2-\omega'^2),\quad
\hbox{or}\quad
\omega^2(u'^2+v'^2+w'^2-\omega'^2).
\]
Here
\[
\begin{aligned}
u&=\alpha p+\beta q+\gamma r,\\
v&=\alpha' p+\beta' q+\gamma' r,\\
w&=\alpha'' p+\beta'' q+\gamma'' r;
\end{aligned}
\]

\newpage
% ===========================================================================
% PDF page 258 -- Book page 236
% ===========================================================================
\setcounter{page}{236}

whence
\[
\begin{aligned}
u'&=\alpha p'+\beta q'+\gamma r',\\
v'&=\alpha' p'+\beta' q'+\gamma' r',\\
w'&=\alpha'' p'+\beta'' q'+\gamma'' r',
\end{aligned}
\]
(the remaining terms vanishing as is well known); and therefore
\[
\begin{aligned}
wv'-v'w&=\alpha(qr'-q'r)+\beta(rp'-r'p)+\gamma(pq'-p'q),\\
wu'-w'u&=\alpha'(qr'-q'r)+\beta'(rp'-r'p)+\gamma'(pq'-p'q),\\
uv'-u'v&=\alpha''(qr'-q'r)+\beta''(rp'-r'p)+\gamma''(pq'-p'q).
\end{aligned}
\]
Hence
\[
\begin{aligned}
wv''-v''w={}&\alpha(qr''-q''r)+\beta(rp''-r''p)+\gamma(pq''-p''q)
             +u'\omega^2-vw\omega',\\
wu''-w''u={}&\alpha'(qr''-q''r)+\beta'(rp''-r''p)+\gamma'(pq''-p''q)
             +v'\omega^2-wu\omega',\\
uv''-u''v={}&\alpha''(qr''-q''r)+\beta''(rp''-r''p)+\gamma''(pq''-p''q)
             +w'\omega^2-uv\omega',
\end{aligned}
\]
and multiplying these by $u'$, $v'$, $w'$, and adding, the required equation is
immediately obtained.

In fact, if $r$ be the distance of a point in the instantaneous axis from the vertex,
and $\rho$, $\sigma$ the radii of curvature of the two cones at that point, then
\[
\frac{r}{\rho}=\frac{\omega^3}{M^2}\Omega,\qquad
\frac{r}{\sigma}=\frac{\omega^3}{M^2}\nabla,
\]
as may be shown without difficulty: and the angular velocity of the instantaneous axis
is given by the equation
\[
\varpi=\frac{M^{3/2}}{\omega^2};
\]
hence the relation between the two angular velocities is
\[
\omega:\varpi=\frac{1}{\rho}-\frac{1}{\sigma}:\frac{1}{r}.
\]

\newpage
% ===========================================================================
% PDF page 259 -- Book page 237 -- Paper [37]
% ===========================================================================
\section*{37.}
\section*{ON THE ROTATION OF A SOLID BODY ROUND A FIXED POINT.}
\addcontentsline{toc}{section}{37. On the Rotation of a Solid Body round a Fixed Point}
\setcounter{page}{237}

\textit{[From the Cambridge and Dublin Mathematical Journal, vol.~I. (1846), pp.~167--173
and 264--274.]}

\bigskip

The difficulty of completing elegantly the solution of this problem, in the case where
no forces act upon the body, arises from the complexity and want of symmetry of the
ordinary formulae for determining the position of one set of rectangular axes with
respect to another set; in consequence of which it has hitherto been considered
necessary to make a particular supposition relative to the position of the fixed axes in
space, viz. that one of them shall be perpendicular to the ``invariable plane'' of the
rotating body. But some formulae for the above purpose, given also by Euler, are
entirely free from these objections. Imagine two sets of axes $Ax$, $Ay$, $Az$, $Ax'$,
$Ay'$, $Az'$. The former set can be made to coincide with the second set, by a rotation
$\theta$ round a certain axis $AR$, inclined to $Ax$, $Ay$, $Az$ at angles $f$, $g$, $h$.
(As usual $f$, $g$, $h$ are the angles $RAx$, $RAy$, $RAz$ considered as positive, and
the rotation is in the same direction as a rotation round $Az$ from $x$ towards $y$.)
This axis may be termed the resultant axis, and the angle $\theta$ the resultant
rotation. The formulae of Euler express the coefficients of the transformation in terms
of the resultant rotation and of the position of the resultant axis, i.e. in terms of
$\theta$ and of the angles $f$, $g$, $h$, whose cosines are connected by the equation
\[
\cos^2 f+\cos^2 g+\cos^2 h=1.
\]
This idea was improved upon by M. Rodrigues (\textit{Liouv.} tom.~v. p.~404), who
introduced the quantities
\[
\tan\tfrac12\theta\cos f,\qquad
\tan\tfrac12\theta\cos g,\qquad
\tan\tfrac12\theta\cos h,
\]
(quantities which will be represented by $\lambda$, $\mu$, $\nu$) by means of which he
expressed the

\newpage
% ===========================================================================
% PDF page 260 -- Book page 238
% ===========================================================================
\setcounter{page}{238}

coefficients as fractions, the numerators of which are very simple rational functions of
the second order of $\lambda$, $\mu$, $\nu$, and which have the common denominator
$(1+\lambda^2+\mu^2+\nu^2)$. These quantities may conveniently be termed the
``coordinates of the resultant rotation,'' and the denominator or the square of the
secant of the semi-angle of resultant rotation will be the ``modulus'' of the rotation.
The elegance of these results led me to apply them to the mechanical question, and I
gave in the \textit{Journal} (vol.~III. p.~224), [6], the differential equations of motion
obtained in terms of $\lambda$, $\mu$, $\nu$: which I integrated as in the common
theory, by supposing one of the fixed axes to be perpendicular to the invariable plane.
Though my attention was again called to the subject, by the connexion of some of these
formulae with Sir William Hamilton's theory of quaternions, no other way of performing
the integration occurred to me. The grand discovery however of Jacobi, of the
possibility of reducing to quadratures the two final differential equations of any
mechanical problem, when the remaining integrals are known, induced me to resume the
problem, and at least attempt to bring it so far as to obtain a differential equation of
the first order between two variables only, the multiplier of which could be obtained
theoretically by Jacobi's discovery. The choice of two new variables to which the
equations of the problem led me, enabled me to effect this with the greatest simplicity;
and the differential equation which I finally obtained, turned out to be integrable
\textit{per se}, so that the laborious process of finding the multiplier became unnecessary.
The new variables $\Omega$, $v$ have the following geometrical interpretations,
$\Omega=k\tan\tfrac12\theta\cos I$, where $k$ is the principal moment, $\theta$ as
before the angle of resultant rotation, and $I$ is the inclination of the resultant axis
to the perpendicular upon the invariable plane, and $v=k^2\cos^2\tfrac12 J$; where, if
we imagine a line $AQ$ having the same position relatively to the axes in fixed space
that the perpendicular upon the invariable plane has to the principal axes of the
rotating body, then $J$ is the inclination of this line to the above perpendicular. To
the choice of these variables I was led by the analysis only. It will be seen that
$p$, $q$, $r$ are functions of $v$ only, while $\lambda$, $\mu$, $\nu$ contain besides the
variable $\Omega$. In obtaining these relations a singular equation
$\Omega^2=\kappa v-k^2$ occurs (equation 13), which may also be written
$1+\tan^2\tfrac12\theta\cos^2 I=\sec^2\tfrac12\theta\cos^2\tfrac12 J$, in which form
the interpretation of the quantities $I$, $J$ has just been given. The equation (17),
it may be remarked, is self-evident: it expresses that the inclination of the resultant
axis to the normal of the invariable plane, is equal to the inclination of the same axis
to the line $AQ$. Now the resultant axis having the same inclination to the axes fixed
in space as it has to the principal axes, and the line $AQ$ the same inclinations to
these fixed axes that the normal to the invariable plane has to the principal axes, the
truth of the proposition becomes manifest. The correspondence in form between the
systems (10) and (14) is also worth remarking. The final results at which I arrive are,
that the time and the arc whose tangent is $\Omega\div k$, are each of them expressible
as the integrals of certain algebraical functions of $v$. The notation throughout is the
same as that made use of in the paper already quoted.

The equations of rotatory motion are
\[
dt=\frac{dp}{P}=\frac{dq}{Q}=\frac{dr}{R}
  =\frac{d\lambda}{\Lambda}=\frac{d\mu}{M}=\frac{d\nu}{N}.\tag{1}
\]

\newpage
% ===========================================================================
% PDF page 261 -- Book page 239
% ===========================================================================
\setcounter{page}{239}

where
\[
\left.
\begin{aligned}
P=\frac{1}{A}\biggl[&(B-C)qr+\frac12\left\{(1+\lambda^2)\frac{dV}{d\lambda}
      +(\lambda\mu+\nu)\frac{dV}{d\mu}
      +(\lambda\nu-\mu)\frac{dV}{d\nu}\right\}\biggr],\\
Q=\frac{1}{B}\biggl[&(C-A)rp+\frac12\left\{(\mu\lambda-\nu)\frac{dV}{d\lambda}
      +(1+\mu^2)\frac{dV}{d\mu}
      +(\mu\nu+\lambda)\frac{dV}{d\nu}\right\}\biggr],\\
R=\frac{1}{C}\biggl[&(A-B)pq+\frac12\left\{(\nu\lambda+\mu)\frac{dV}{d\lambda}
      +(\mu\nu-\lambda)\frac{dV}{d\mu}
      +(1+\nu^2)\frac{dV}{d\nu}\right\}\biggr],
\end{aligned}
\right\}\tag{2}
\]
\[
\left.
\begin{aligned}
\Lambda&=\frac12\{(1+\lambda^2)p+(\lambda\mu-\nu)q+(\lambda\nu+\mu)r\},\\
M&=\frac12\{(\mu\lambda+\nu)p+(1+\mu^2)q+(\mu\nu-\lambda)r\},\\
N&=\frac12\{(\nu\lambda-\mu)p+(\mu\nu+\lambda)q+(1+\nu^2)r\},
\end{aligned}
\right\}\tag{3}
\]
[whence also
\[
\lambda\Lambda+\mu M+\nu N=\tfrac12\kappa(\lambda p+\mu q+\nu r)\tag{3 bis}
\]
].

In the case where the forces vanish, the first three equations become simply
\[
\left.
\begin{aligned}
P&=\frac{1}{A}(B-C)qr,\\
Q&=\frac{1}{B}(C-A)rp,\\
R&=\frac{1}{C}(A-B)pq,
\end{aligned}
\right\}\tag{4}
\]
and here the usual four integrals of the system are
\[
Ap^2+Bq^2+Cr^2=h,\tag{5}
\]
\[
\left.
\begin{aligned}
Ap(1+\lambda^2-\mu^2-\nu^2)+2Bq(\lambda\mu-\nu)+2Cr(\nu\lambda+\mu)
 &=a(1+\lambda^2+\mu^2+\nu^2),\\
2Ap(\lambda\mu+\nu)+Bq(1+\mu^2-\nu^2-\lambda^2)+2Cr(\mu\nu-\lambda)
 &=b(1+\lambda^2+\mu^2+\nu^2),\\
2Ap(\nu\lambda-\mu)+2Bq(\mu\nu+\lambda)+Cr(1+\nu^2-\lambda^2-\mu^2)
 &=c(1+\lambda^2+\mu^2+\nu^2),
\end{aligned}
\right\}\tag{6}
\]
or as they may also be written,
\[
\left.
\begin{aligned}
a(1+\lambda^2-\mu^2-\nu^2)+2b(\lambda\mu+\nu)+2c(\nu\lambda-\mu)
 &=Ap(1+\lambda^2+\mu^2+\nu^2),\\
2a(\lambda\mu-\nu)+b(1+\mu^2-\nu^2-\lambda^2)+2c(\mu\nu+\lambda)
 &=Bq(1+\lambda^2+\mu^2+\nu^2),\\
2a(\nu\lambda+\mu)+2b(\mu\nu-\lambda)+c(1+\nu^2-\lambda^2-\mu^2)
 &=Cr(1+\lambda^2+\mu^2+\nu^2).
\end{aligned}
\right\}\tag{6 bis}
\]
to which we may add,
\[
A^2p^2+B^2q^2+C^2r^2=k^2,\tag{7}
\]
where
\[
k^2=a^2+b^2+c^2.\tag{8}
\]
Introducing the quantities $\kappa$, $\Omega$, (the former of which has been already
made use of) given by the equations
\[
\left.
\begin{aligned}
\kappa&=1+\lambda^2+\mu^2+\nu^2,\\
\Omega&=\lambda Ap+\mu Bq+\nu Cr,
\end{aligned}
\right\}\tag{9}
\]

\newpage
% ===========================================================================
% PDF page 262 -- Book page 240
% ===========================================================================
\setcounter{page}{240}

The equations (6) may be written under the form
\[
\left.
\begin{aligned}
2\lambda\Omega+2\mu Cr-2\nu Bq&=\kappa(Ap+a)-2Ap,\\
-2\lambda Cr+2\mu\Omega+2\nu Ap&=\kappa(Bq+b)-2Bq,\\
2\lambda Bq-2\mu Ap+2\nu\Omega&=\kappa(Cr+c)-2Cr,
\end{aligned}
\right\}\tag{10}
\]
whence also, multiplying by $Ap$, $Bq$, $Cr$, and adding,
\[
2\Omega^2=\kappa\{k^2+(Apa+Bqb+Crc)\}-2k^2,\tag{11}
\]
or writing
\[
k^2+(Apa+Bqb+Crc)=2v,\tag{12}
\]
this becomes
\[
\Omega^2=\kappa v-k^2:\tag{13}
\]
an equation, the geometrical interpretation of which has already been given.

From the equations (10) we deduce the inverse system
\[
\left.
\begin{aligned}
a\Omega-bCr+cBq&=2\lambda v-\Omega Ap,\\
aCr+b\Omega-cAp&=2\mu v-\Omega Bq,\\
-aBq+bAp+c\Omega&=2\nu v-\Omega Cr,
\end{aligned}
\right\}\tag{14}
\]
which are easily verified by multiplying by $\Omega$, $Cr$, $-Bq$; or by $-Cr$,
$\Omega$, $Ap$; or $Bq$, $-Ap$, $\Omega$: adding and reducing, by which means the
equations (10) are re-obtained. Hence also if for shortness
\[
\left.
\begin{aligned}
\Phi&=ap+bq+cr,\\
\nabla&=aqr(B-C)+brp(C-A)+cpq(A-B),
\end{aligned}
\right\}\tag{15}
\]
we have, multiplying by $p$, $q$, $r$, and adding,
\[
\Omega\Phi-\nabla=2v(\lambda p+\mu q+\nu r)-\Omega h.\tag{16}
\]
To these may be added the equation
\[
\Omega=a\lambda+b\mu+c\nu,\tag{17}
\]
which follows immediately from either of the systems (10) or (14).

We may also put the equations (10) under this other form,
\[
\left.
\begin{aligned}
2\lambda\Omega-2\mu c+2\nu b&=\kappa(Ap+a)-2a,\\
2\lambda c+2\mu\Omega-2\nu a&=\kappa(Bq+b)-2b,\\
-2\lambda b+2\mu a+2\nu\Omega&=\kappa(Cr+c)-2c,
\end{aligned}
\right\}\tag{10 bis}
\]
It may be remarked now, that $p$, $q$, $r$ are functions of $v$; since we have to
determine these quantities, the three equations
\[
\left.
\begin{aligned}
Ap^2+Bq^2+Cr^2&=h,\\
A^2p^2+B^2q^2+C^2r^2&=k^2,\\
Apa+Bqb+Crc&=2v-k^2,
\end{aligned}
\right\}\tag{18}
\]

\end{document}
